%% parts/p07/cut-elimination-and-interpolation-supplement.tex
% Necessity: the exact inherited cut-elimination chapter states Craig and Beth
% but does not contain the separately attributed midsequent, Herbrand, or
% Maehara apparatus required by the open depth and new-content registers.
% The partition proof is followed by direct relational cut-free completeness,
% graph relationalization for functions and equality, and the Beth argument.
% The original Craig and Beth labels remain their public destinations.
% Use: the Part VII driver routes it after the six inherited cut-chapter wrappers;
% it is not represented as a verbatim restored slice.

\section{Midsequents and the Herbrand route}

\begin{hbdefinition}{def:p07-regular-cut-free-fol-derivation}{Regular cut-free derivation}
\par\noindent
For this section, a \emph{regular derivation} is a
finite \(\mathsf{LK}_{\mathrm d}\)-derivation that contains no Cut, Mix,
equality, or induction inference and assigns distinct eigenvariables to
distinct quantifier inferences, each fresh for the conclusion required by its
rule.  This is a metalinguistic classification of derivation trees, not a new
calculus.  Regularity is needed for the quantifier permutations in the next
theorem and differs from cut-freeness by its global eigenvariable condition.
\index{regular derivation}
\index{midsequent theorem}
\end{hbdefinition}

\begin{hbtheorem}{thm:p07-midsequent}{cited=Takeuti1987}{Midsequent theorem, equality-free prenex form}
\par\noindent
If a prenex endsequent has a regular equality-free cut-free derivation on
\(\mathsf{LK}_{\mathrm d}\), that derivation can be permuted so that no
propositional logical inference occurs below a quantifier inference, with
tree roots drawn at the bottom.  A quantifier-free sequent then separates the
upper propositional part from the lower quantifier part.
\end{hbtheorem}

\begin{hbtheorem}{thm:p07-herbrand-midsequent}{cited=Takeuti1987}{Herbrand route from a midsequent}
\par\noindent
Let
\[
 \exists x_1\cdots\exists x_k\,\varphi(x_1,\ldots,x_k)
\]
be a valid equality-free prenex sentence with quantifier-free matrix, in a
signature containing a ground term or after adding one fresh constant.  A cut-free
derivation followed by the midsequent permutation exposes finitely many
ground tuples \((t_1^j,\ldots,t_k^j)\), \(1\leq j\leq n\), for which
\[
 \varphi(t_1^1,\ldots,t_k^1)\lor\cdots\lor
 \varphi(t_1^n,\ldots,t_k^n)
\]
is propositionally valid.
\end{hbtheorem}

This is a proof-theoretic route to the citation-only validity form in
Theorem~\ref{thm:p02-herbrand-validity}; it neither upgrades that theorem nor
replaces its finite-clause unsatisfiability form or the epsilon-calculus
route.  An equality-bearing Herbrand theorem would require equality instances
or an equality calculus and is not asserted.

\section{Maehara interpolation and definability}

\begin{hbdefinition}{def:p07-partition-interpolant}{Partition and partition interpolant}
For the local partition lemma, fix a \emph{purely relational signature}:
equality, function symbols, and constant symbols are absent, so its terms are
variables.  This is the fragment used by the quantifier induction in
Annex~\ref{annex:p7.interp.1}.  For general first-order signatures,
Annex~\ref{annex:p7.interp.3} represents functions by relations and equality
by an equivalence relation, applies this partition lemma, and translates
the resulting formula back into the common signature.

Let an equality-free sequent in that signature be written
\(\Gamma_1,\Gamma_2\Rightarrow\Delta_1,\Delta_2\).  The ordered four-tuple of
displayed subcontexts is a \emph{partition} of the sequent.  A
\emph{partition interpolant} is a formula \(\chi\) such that
\(\Gamma_1\Rightarrow\Delta_1,\chi\) and
\(\chi,\Gamma_2\Rightarrow\Delta_2\) are derivable in
\(\mathsf{LK}_{\mathrm d}\), and every nonlogical
predicate symbol in \(\chi\) occurs both in a formula
of \(\Gamma_1\cup\Delta_1\) and in a formula of
\(\Gamma_2\cup\Delta_2\).  In addition,
\[
 \FV(\chi)\subseteq
 \FV(\Gamma_1\cup\Delta_1)\cap
 \FV(\Gamma_2\cup\Delta_2).
\]
Here \(\bot\) is logical and equality is absent.  Thus a partition whose four
contexts contain only sentences has a sentence interpolant.  The free-variable
clause is required because quantifier-rule premises may contain parameters even
when the final sequent contains none.

The term is necessary because an arbitrary intermediate consequence lacks
the common-vocabulary and free-variable clauses; it names a formula relative to a partition,
not a new consequence relation.  That clause is used to obtain Craig's
vocabulary restriction below.
\index{partition interpolant}
\end{hbdefinition}

\begin{hbmethod}{def:p07-maehara-method}{Maehara's partition method}
Induct on a supplied cut-free equality-free
\(\mathsf{LK}_{\mathrm d}\)-derivation together with a partition of its
endsequent.  Construct an interpolant at initial and \(\bot L\) leaves;
propagate it through weakening and contraction; combine premise interpolants
separately for every placement of each propositional principal formula; and
  in the four quantifier cases quantify away eigenvariables and instantiation
  variables not common to the conclusion partition while preserving the two
  required derivations.  Applied to a
cut-free derivation of \(\varphi\Rightarrow\psi\) with partition
\((\varphi;\varnothing\mid\varnothing;\psi)\), the output has the common
nonlogical vocabulary needed for Craig interpolation.  The annex carries the
complete initial, structural, propositional, and quantifier induction.
\end{hbmethod}

\begin{hblemma}{thm:p07-maehara-partition}{proved=annex:p7.interp.1}{Maehara's partition lemma}
Every partition of the endsequent of a cut-free equality-free derivation in
the purely relational restriction of \(\mathsf{LK}_{\mathrm d}\) has a partition interpolant in the
sense of Definition~\ref{def:p07-partition-interpolant}.  In particular, if
the four endsequent contexts contain only sentences, the interpolant is a
sentence.
\end{hblemma}

\begin{hbtheorem}{thm:p07-relational-cut-free-completeness}{proved=annex:p7.interp.2}{Relational soundness and cut-free completeness}
Let \(\Gamma\Rightarrow\Delta\) be a finite sequent in an equality-free
purely relational signature.  It is valid in all nonempty classical
structures, under every assignment, if and only if it has a derivation
in \(\mathsf{LK}_{\mathrm d}\) containing neither Cut nor Mix.
Every derivation in this fragment, including those using Cut or Mix,
has a valid endsequent.  Empty contexts and nullary predicates are included.
\end{hbtheorem}

The proof constructs a countermodel from a fair sequence of unprovable
sequents.  Its domain is the set of variables used as parameters, and its
relations are determined by antecedent atoms.  Together with Maehara's
lemma, this gives semantic interpolation in the relational fragment.
The search and model method is developed in
\cite[\S5.4]{Gallier2015}; the annex verifies the rules of the
particular calculus used here.

\begin{center}
\begin{tabularx}{\textwidth}{@{}p{.23\textwidth} X X@{}}
\toprule
Derivation case & Interpolant construction & Vocabulary/decrease check \\
\midrule
\(A\Rightarrow A\) initial & \(\bot,\top,A,\neg A\) when the two occurrences lie respectively left/left, right/right, left/right, right/left of the partition & nonlogical symbols occur on both sides only in the cross cases \\
\(\bot L\) initial & \(\bot\) if the principal occurrence is in the first part, otherwise \(\top\) & both choices are logical constants \\
exchange, weakening, contraction & retain the premise interpolant & no new symbol; the required sequents follow by the same structural rule \\
one-premise propositional rule & retain the premise interpolant & apply the same logical rule on the partition side containing the principal formula \\
two-premise propositional rule & use \(I_0\lor I_1\) when the principal formula is in the first part and \(I_0\land I_1\) when it is in the second & \(\lor L/\land R\) combine the opposite-side derivations; only common symbols from the premises remain \\
quantifier rule without an eigenparameter & retain \(I\) if its variables remain common; otherwise use \(\forall xI(x)\) for a principal formula in the first part and \(\exists xI(x)\) for one in the second & the quantified variable is precisely the instantiation variable that ceases to be common \\
eigenparameter rule & if fresh \(a\) occurs in the constructed \(I(a)\), use \(\exists xI(x)\) for a principal formula in the first part and \(\forall xI(x)\) for one in the second & the quantifier removes \(a\), and the dual quantifier rule reconstructs both partition sequents \\
\bottomrule
\end{tabularx}
\end{center}

\paragraph{Proof location for the partition lemma.}
The table is a reader map, not a substitute for the induction.  The complete
construction, including all four placements of an initial formula, both
placements of every principal formula, the three branching rules, and both
eigenvariable rules, is Annex~\ref{annex:p7.interp.1}
\cite{Maehara1960}.

\paragraph{Functions, equality, and common vocabulary.}
Only finitely many symbols occur in the two formulas of an interpolation
problem.  Replace each function by its graph relation and impose totality
and uniqueness; replace equality by an equivalence relation and impose
congruence for every predicate and graph.  A structure satisfying these
axioms has a quotient that interprets the original functions and equality.
The quotient truth lemma transfers the original valid implication to the
relational fragment.

The axioms for symbols occurring only on the first side accompany its
antecedent.  The axioms for symbols occurring only on the second side
accompany the antecedent of its conclusion implication.  Consequently the
relational interpolant can mention only common graphs and predicates,
together with the equality relation.  Replacing those graphs by equations
produces the required first-order interpolant.
Annex~\ref{annex:p7.interp.3} gives the axioms, both truth inductions, the
vocabulary argument, and the elimination of fresh constants used for free
variables.

For example, let
\[
\begin{aligned}
 \varphi&=P(c)\land\forall x\,(P(x)\to Q(f(x))),\\
 \psi&=\bigl(\forall z\,(Q(z)\to R(z))\bigr)\to\exists z\,R(z).
\end{aligned}
\]
Then \(\theta=\exists z\,Q(z)\) is an interpolant.  The first implication
uses the witness \(f(c)\).  For the second, a witness to \(Q\) also
witnesses \(R\) under the antecedent of \(\psi\).
Only \(Q\) is shared; \(P,f,c\), and \(R\) do not occur in \(\theta\).

\paragraph{From implicit to explicit definition.}
The Beth proof uses two copies of the distinguished relation while retaining
one shared reduct.  Compactness extracts finite subtheories, and Craig
interpolation removes both copies of the relation from the intermediate
formula.  Renaming the second copy back gives the two implications of an
explicit definition.  Annex~\ref{annex:p7.interp.4} carries the full argument,
including an infinite starting theory, arity zero, and fresh-parameter
elimination.

\paragraph{Example: a relation constrained to choose right inverses.}
Use a binary function \(\cdot\), a constant \(e\), and a binary predicate
\(R\).  Let \(T\) contain
\[
\begin{gathered}
 \forall x\,y\,z\,((x\cdot y)\cdot z=x\cdot(y\cdot z)),\\
 \forall x\,(e\cdot x=x\land x\cdot e=x),\\
 \forall x\,\exists y\,R(x,y),\qquad
 \forall x\,y\,(R(x,y)\to x\cdot y=e).
\end{gathered}
\]
The relation is total and contains only right-inverse pairs.
Every element therefore has a right inverse.
If \(a\cdot b=e\), choose \(c\) with \(b\cdot c=e\).  Then
\[
 a=a\cdot e=a\cdot(b\cdot c)
   =(a\cdot b)\cdot c=e\cdot c=c,
\]
so \(b\cdot a=e\).
If also \(a\cdot d=e\), then
\[
 d=e\cdot d=(b\cdot a)\cdot d
   =b\cdot(a\cdot d)=b\cdot e=b.
\]
Thus each element has a unique right inverse.  Totality forces \(R\)
to contain it, and the last axiom excludes every other pair.
Consequently \(T\) implicitly defines \(R\), with the explicit definition
\[
 T\models\forall x\,y\,
             (R(x,y)\leftrightarrow x\cdot y=e).
\]
The defining formula uses only the function, the constant, and logical
equality.  The four displayed sentences themselves form a finite support.
